% Zone „Das kennst du schon“ – aus bank/kreis/zone.jsonl (zusammenbau v0.1)
\einheitenkopf[zone][Kennst du schon]{Das kennst du schon}
\setcounter{aufgabe}{0}

%% TODO Titel: Ich-kann-Satz fehlt in der Bank (Platzhalter: Kettenname) – Nr. 1
%% TODO Anweisung: ein Satz über der ersten Teilaufgabe fehlt in der Bank – Nr. 1
\begin{aufgabe}{Fläche und Umfang unterscheiden; Einheiten cm, cm², m, m²}
\begin{teile}
\teil Rechteck $5$ cm lang, $3$ cm breit – wie groß ist der Umfang? \leerfeld[cm]
\teil Rechteck $4$ m lang, $1{,}5$ m breit – Umfang und Fläche? Umfang = \leerfeld[m] , Fläche = \leerfeld[m²]
\end{teile}
\end{aufgabe}

%% TODO Titel: Ich-kann-Satz fehlt in der Bank (Platzhalter: Kettenname) – Nr. 2
%% TODO Anweisung: ein Satz über der ersten Teilaufgabe fehlt in der Bank – Nr. 2
\begin{aufgabe}{Mit Dezimalzahlen multiplizieren und dividieren, Taschenrechner mit π, Ergebnis runden (Näherungswert mal Kommazahl)}
\begin{teile}
\teil $0{,}5 \cdot 8$ – wie viel? \leerfeld
\teil $2{,}73 \cdot 5{,}4$ – wie viel? Runde auf eine Stelle. ≈ \leerfeld
\end{teile}
\end{aufgabe}

%% TODO Titel: Ich-kann-Satz fehlt in der Bank (Platzhalter: Kettenname) – Nr. 3
%% TODO Anweisung: ein Satz über der ersten Teilaufgabe fehlt in der Bank – Nr. 3
\begin{aufgabe}{Formel nach einer Größe umstellen (u = π · d → d = u : π)}
\begin{teile}
\teil Ein Rechteck hat die Fläche $24$ cm² und ist $6$ cm lang – wie breit ist es? \leerfeld[cm]
\teil Für ein Rechteck gilt $A = a \cdot b$ ($A$ Fläche, $a$ Länge, $b$ Breite). Stelle nach $b$ um und berechne $b$ für $A = 17{,}5$ m² und $a = 3{,}5$ m. b = \leerfeld[m]
\end{teile}
\end{aufgabe}

%% TODO Titel: Ich-kann-Satz fehlt in der Bank (Platzhalter: Kettenname) – Nr. 4
%% TODO Anweisung: ein Satz über der ersten Teilaufgabe fehlt in der Bank – Nr. 4
\begin{aufgabe}{Quadrieren und Wurzelziehen (4,5²; √20,25)}
\begin{teile}
\teil $7^2$ – wie viel? \leerfeld
\teil $\sqrt{38}$ – wie viel? Runde auf zwei Stellen. ≈ \leerfeld
\end{teile}
\end{aufgabe}

%% TODO Titel: Ich-kann-Satz fehlt in der Bank (Platzhalter: Kettenname) – Nr. 5
%% TODO Anweisung: ein Satz über der ersten Teilaufgabe fehlt in der Bank – Nr. 5
\begin{aufgabe}{Winkel messen und zeichnen, Vollwinkel}
\begin{teile}
\teil Ein halber rechter Winkel – wie viel Grad? \leerfeld °
\teil Miss den Winkel $\alpha$ mit dem Geodreieck.

\winkel[4]{65}{\alpha}
α = \leerfeld °
\end{teile}
\end{aufgabe}

%% TODO Titel: Ich-kann-Satz fehlt in der Bank (Platzhalter: Kettenname) – Nr. 6
%% TODO Anweisung: ein Satz über der ersten Teilaufgabe fehlt in der Bank – Nr. 6
\begin{aufgabe}{Anteil als Bruch und Prozentsatz bilden (Mittelpunktswinkel zum Vollwinkel als Prozent)}
\begin{teile}
\teil $6$ von $24$ – welcher Bruchteil? Kürze. \leerfeld
\teil $17$ von $45$ – wie viel Prozent? Runde auf eine Stelle. ≈ \leerfeld[\%]
\end{teile}
\end{aufgabe}

%% TODO Titel: Ich-kann-Satz fehlt in der Bank (Platzhalter: Kettenname) – Nr. 7
%% TODO Anweisung: ein Satz über der ersten Teilaufgabe fehlt in der Bank – Nr. 7
\begin{aufgabe}{Fläche und Umfang unterscheiden; Einheiten cm, cm², m, m² – Fehler finden}
\begin{teile}
\teil Ein Beet ist $6$ m lang und $4$ m breit. Es bekommt rundherum einen Zaun. Paul rechnet: \rechnung{6 \cdot 4 &= 24 \\ \text{Zaun} &= 24 \text{ m}} Wo ist der Fehler?
\end{teile}
\end{aufgabe}

%% TODO Titel: Ich-kann-Satz fehlt in der Bank (Platzhalter: Kettenname) – Nr. 8
%% TODO Anweisung: ein Satz über der ersten Teilaufgabe fehlt in der Bank – Nr. 8
\begin{aufgabe}{Fläche und Umfang unterscheiden; Einheiten cm, cm², m, m² – selbst rechnen}
\begin{teile}
\teil Ein Spielfeld ist $9$ m lang und $5$ m breit. Rundherum kommt ein Absperrband. Wie viel Band braucht man? \leerfeld[m]
\end{teile}
\end{aufgabe}
